\section{Détermination de la constante et covariance de ses réalisations}
\label{act21:centralidad-constante}

La constante holographique de couplage arithmético--géométrique réunit une
détermination structurelle et une capacité de restitution. La première fixe
le registre ordonné qui est compatible avec les incidences et avec la
chronologie de l’état engendré. La seconde établit comment récupérer ce
registre à partir de ses représentations et comment reconnaître les opérations qui
agissent sur lui. La composition des deux propriétés donne un contenu précis
à sa fonction de couplage.

La détermination commence dans les régions produites par APP--TRIT--TPK. Leurs bandes de six trits conservent la position de chaque symbole, la charge, la retenue et la profondeur. L'horloge entière compare les échelles ternaire et décimale ; sa première avancée nulle, suivie du retour nonadique, détermine la lecture phase--charge du descripteur. Le répertoire terminal et ses incidences sont conservés avec cette provenance. Les ordonnancements de ce répertoire produisent le domaine fini sur lequel agit la sélection conjointe de la section~\ref{act21:chap:despliegue-selector-k}.

La condition exceptionnelle restreint ce domaine aux registres dont la face
de support coïncide avec la face extraite des mots régionaux. La
condition temporelle compare, dans un même ordre, les lectures
comparative et affine des deux tronçons terminaux. Ainsi, la position
d’une coordonnée et la relation entre deux instants participent à la même
décision. Le résultat exact est
\[
 19\,446\longrightarrow79\longrightarrow1.
\]
L’unique registre survivant détermine l’entier
\[
 N_K=234543140729659824621058914794146601,
 \qquad
 \kappa_{\rm ag}=\frac{N_K}{10^{36}-1}.
\]
L’expression rationnelle conserve l’ordre des douze triades dans le système de coordonnées
de phase spécifié. L’unicité se réfère au domaine terminal
construit, avec ses deux conditions explicites et ses données d’orientation.
Ainsi est fixé ce qui est sélectionné, au moyen de quelles relations et dans
quel ensemble de possibilités.

Le recensement des trois terminaisons régionales conserve une autre fonction.
Ses \(196\cdot197\cdot196\) éléments sont des triplets de mots de six
trits compatibles avec les cylindres déclarés. Le profil de Gram, les
poids, les distances et les marges orientées déterminent la terminaison
visible. Le recensement de \(19\,446\) éléments examine des registres ordonnés à
douze coordonnées. Les deux reçoivent de l’information de la même histoire
enrichie ; leurs types et leurs filtres permettent de suivre comment ils sont reliés
sans confondre un triplet visible avec le registre complet.

Une fois \(K\) sélectionné, la transformation de Hadamard et ses lecteurs établissent des équivalences entre représentations de ce même registre. La sélection apporte la détermination parmi les candidats ; l'inversion apporte la restitution à partir d'une représentation. La loi bilatérale fournit ensuite l'identité de conservation pour les isométries déclarées, et l'archive conserve l'information qui permet d'inverser des lectures successives. Chaque résultat transmet au suivant un objet avec domaine, orientation et normalisation définis.

Cette hiérarchie précise également le rôle de \(K\) dans la loi de
conservation. L’identité bilatérale est démontrée pour deux isométries
ayant un domaine et un codomaine communs. Dans les réalisations de l’article, le
registre sélectionné fixe les lectures arithmétiques et incidencielles,
sa direction centrée et les opérations de reconstruction qui se
composent avec cette identité. La généralité du bilan et la
détermination du registre sont des propriétés complémentaires : la première
décrit une loi opératoire ; la seconde identifie la structure sur
laquelle sont calculées les réalisations exposées ici.

La clôture de la constante de structure fine utilise les coordonnées
régionales déjà engendrées et le registre sélectionné, en conservant les
retenues de l’opération dodécaphasique. Sa représentation analytique
développe la même précoordonnée avec la récurrence complète. Par
l’autre projection, les incidences conservées sont transportées aux
réalisations exceptionnelles. La provenance commune permet de relier
valeur arithmétique, incidence et mémoire sans convertir leurs différents
lecteurs en une grandeur unique.

Les réalisations géométriques ultérieures élargissent cette fonction. La direction centrée de \(K\) détermine le flux diagonal de déterminant un, préserve les signes de Plücker et le covolume des réseaux spécifiés ; l'archive permet de la récupérer et de contrôler ses erreurs de lecture. La transformation thêta--Mellin conserve la structure polaire dans le domaine démontré, tandis que la différence entière enregistre la variation spectrale. La sélection temporelle est ainsi liée, au moyen des opérations déjà exposées, à une structure qui admet évaluation, transport et récupération.

L’apport conjoint de l’article s’organise donc en une
suite démonstrative définie : génération de l’état, sélection de la
constante, équivalence de ses représentations, conservation opératoire
et restitution stable. Cette suite détermine la position de chaque
théorème et maintient visibles les conditions qu’il reçoit et les relations
qu’il transmet.
